%% ==================== RELATED WORK ====================
\section{Related Work}
\label{sec:related}

Three literatures bear on this work: visual-inertial SLAM, which supplies the tracking core; brain-inspired spatial cognition, which supplies the state representation; and learning-based visual processing, which supplies much of the front-end. Each is summarised below with an explicit statement of where NeuroLocMap departs from it, and Table~\ref{tab:comparison_bioslam} collects the architectural differences.

\subsection{Visual SLAM and Sensor Fusion}
\label{sec:related_slam}

Visual SLAM has developed along two lines. Feature-based systems such as ORB-SLAM~\cite{mur2015orb} extract sparse keypoints and reach real-time performance in textured scenes; direct systems such as LSD-SLAM~\cite{engel2014lsd} and DSO~\cite{engel2018direct} operate on pixel intensities without feature extraction and recover semi-dense structure. Place recognition and loop closure have their own tradition, from the probabilistic bag-of-words model of FAB-MAP~\cite{Cummins2008} to the learned descriptors of NetVLAD~\cite{Arandjelovic2016}. Two constraints recur across this literature. Feature-based pipelines require texture, typically hundreds of keypoints per frame, and degrade in its absence; direct pipelines require photometric stability and tolerate neither motion blur nor uncalibrated exposure. In both cases the dominant cost is attached to the visual measurements themselves, whether through extraction and matching or through dense photometric optimisation.

Fusing an IMU with the camera removes part of that dependence, since inertial data is unaffected by illumination and supplies the high-frequency motion that vision misses. MSCKF~\cite{Mourikis2007} and VINS-Mono~\cite{Qin2018} exemplify the tightly coupled formulation, in which feature observations enter the state and the estimator is solved as one optimisation. Complementary filtering takes the opposite view: Mahony et al.~\cite{Mahony2008} and Madgwick et al.~\cite{Madgwick2011} divide the labour by frequency band, letting vision carry the low-frequency estimate and the IMU the high-frequency dynamics. Tightly coupled estimators reach high accuracy but require careful bias and scale initialisation and Jacobians over every feature state; the complementary view yields a smaller, more transparent core. NeuroLocMap takes the second route and instantiates it as a fixed-dimension (10-state) loosely coupled EKF, so that the per-update cost does not grow with the feature count, while the low/high-frequency division mirrors the vestibular--visual integration described in the vestibular literature~\cite{Angelaki2004,Cullen2012}.

On the mapping side, Gaussian-splatting systems have recently improved reconstruction within this journal: ReSemGS-SLAM enforces semantic consistency across views~\cite{zhao2026resemgs}, and deblur-aware splatting extends the same representation to degraded imagery~\cite{qin2025deblur}. Both optimise a photometric map and neither estimates monocular scale, which places them in a different problem class from the topological map used here; they are cited because they define the current state of dense mapping, not because they are direct competitors.

\subsection{Brain-Inspired Spatial Cognition}
\label{sec:related_bio}

Biological navigation rests on a small set of spatially selective cell types. Place cells in the hippocampus fire at specific locations and were the first evidence of an internal map~\cite{OKeefe1971}; grid cells in the entorhinal cortex carry periodic hexagonal firing fields that support path integration~\cite{Hafting2005}; head-direction cells encode facing as an internal compass~\cite{Taube1990}. Moser et al.~\cite{Moser2017} review how these populations cooperate, and work on bats established that the same principles extend to volume, with three-dimensional place and grid cells on a face-centred cubic lattice and three-dimensional head-direction cells~\cite{finkelstein2015three,finkelstein2016three}.

Robot implementations of this architecture begin with RatSLAM~\cite{milford2004ratslam,milford2008mapping}, which couples continuous attractor networks for pose cells to an experience map for loop closure and demonstrated suburb-scale operation over 66~km. BatSLAM~\cite{steckel2013batslam} transferred the same scheme to sonar, and NeuroSLAM~\cite{Yu2019} extended pose encoding to three dimensions with CNN-based visual processing. Two recent systems are closest to the present work. ORB-NeuromSLAM couples ORB features with a neuromorphic processing chain for three-dimensional mapping~\cite{shen2024orbneuro}, and Hybrid-NeuroSLAM combines visual and inertial cues through neurobiologically inspired modules at large scale~\cite{lu2025hybridneuro}. Both estimate a full 6-DoF pose and therefore carry roll and pitch in the state; NeuroLocMap commits to 4-DoF instead, which removes those states and the observability problems attached to them, at the price of the planar-motion assumption discussed in Section~\ref{sec:intro}.

Three gaps remain in this line. Most frameworks are planar or encode position without heading, so they lack the 4-DoF state that road and indoor platforms require; visual front-ends rely on intensity templates whose count saturates in complex scenes, RatSLAM yielding on the order of five; and inertial fusion is either absent or ad hoc, so the frequency-domain complementarity that biology exploits is not realised in the estimator. The present system addresses the three in turn: grid cells on a multi-dimensional continuous attractor together with multilayered head-direction cells for the 4-DoF state, a dual-stream front-end that raises the template count from five to 964, and complementary filtering for vestibular--visual fusion.

\subsection{Learning-Based Visual Processing}
\label{sec:related_learning}

Deep networks have changed what a visual front-end can supply. The dual-stream account of vision, ventral for identity and dorsal for motion, is well established biologically~\cite{Goodale1992,Mishkin1983} and has been used to organise recognition architectures~\cite{DiCarlo2012,Duffy1998}; Transformer models~\cite{Vaswani2017} added long-range dependency modelling through self-attention. The deployment record of these methods is more mixed than their benchmark record. Place-recognition embeddings of the kind reviewed above~\cite{Arandjelovic2016} are trained on datasets of city-scale coverage and lose accuracy on unseen domains unless retrained; full self-attention scales quadratically with sequence length, $\mathcal{O}(N^2)$, which rules out its use inside a per-frame budget; and learned features carry no explicit spatial structure, so they cannot be read as a map or coupled to metric constraints without a further layer of interpretation.

NeuroLocMap therefore adopts a narrower contract: everything learned is trained once in simulation and deployed without per-domain retraining. The front-end follows the ventral/dorsal division with a lightweight modulation that captures global--local interaction statistically instead of through full attention matrices, and the loop-closure channel uses a learned embedding whose transfer across domains is measured rather than assumed. The consequence is that template discrimination is set by the front-end rather than by a learned metric, which is what allows the same pipeline to run at 39--76~FPS on a CPU.

\subsection{Comparative Analysis}
\label{sec:related_compare}

Table~\ref{tab:comparison_bioslam} compares representative bio-inspired frameworks along the dimensions that determine what a system can be used for: degrees of freedom, spatial cell types, fusion mechanism, visual processing, loop closure and temporal modelling. Four differences recur. Representation is planar in RatSLAM and BatSLAM, positional without heading in NeuroSLAM, and fully 6-DoF in the two recent systems; 4-DoF is the middle ground this work occupies. Visual processing ranges from one-dimensional scanlines to ORB features, with template diversity as the practical consequence. Fusion ranges from absent to tightly coupled, with the per-update cost following the choice. And nothing in the earlier systems couples attention-style feature enhancement to explicit spatial cells. NeuroLocMap addresses each of the four, and the comparison is architectural: it claims a smaller, more interpretable state at comparable error, not accuracy superiority over 6-DoF systems evaluated on other benchmarks.

\begin{table*}[!t]
       \centering
       \scriptsize
       \caption{Comparison of bio-inspired navigation frameworks}
       \label{tab:comparison_bioslam}
       \setlength{\tabcolsep}{3pt}
       \begin{threeparttable}
              \begin{tabular}{@{}l@{\extracolsep{\fill}}cccccc@{}}
                     \toprule
                     Framework & DoF & Spatial Cells & Fusion & Visual Proc. & Loop Closure & Temp. Model \\
                     \midrule
                     RatSLAM~\cite{milford2008mapping} & 2D & PC+HDC & Vis. & Scanline & Exp. & - \\
                     BatSLAM~\cite{steckel2013batslam} & 2D & PC+HDC & Sonar & Echo & Exp. & - \\
                     NeuroSLAM~\cite{Yu2019} & 3D & 3D-GC & Vis. & CNN & Exp. & LSTM \\
                     \midrule
                     \textbf{NeuroLocMap (Ours)} & \textbf{4D} & \textbf{3D-GC+M-HDC} & \textbf{Comp.} & \textbf{Dual+HART} & \textbf{Exp.} & \textbf{Trans.} \\
                     \bottomrule
              \end{tabular}
              \begin{tablenotes}
                     \scriptsize
                     \item \textbf{Abbreviations}: PC: Place Cells; HDC: Head Direction Cells; GC: Grid Cells; M-HDC: Multilayered HDC; Exp.: Experience map; Comp.: Complementary filtering; Trans.: Transformer-inspired temporal modeling; Vis.: Visual only; Dual+HART: Dual-stream (ventral/dorsal) visual processing with HART-inspired spatial attention.
                     \item \textbf{Note}: Our feature enhancement uses global statistical modulation ($\mathcal{O}(N)$ complexity) to approximate self-attention effects, rather than full Query-Key-Value operations ($\mathcal{O}(N^2)$), enabling real-time performance without training.
              \end{tablenotes}
       \end{threeparttable}
\end{table*}
